% =====================================================================
%  FICHE 8 — THÈME 4 — NIVEAU DIFFICILE — ÉNONCÉS
% =====================================================================
\documentclass[11pt,a4paper]{article}
\usepackage[utf8]{inputenc}
\usepackage[T1]{fontenc}
\usepackage{lmodern}
\usepackage[french]{babel}
\usepackage{amsmath,amssymb,mathtools}
\usepackage{icomma}
\usepackage{siunitx}
\sisetup{output-decimal-marker={,},per-mode=power,group-separator={\,},group-minimum-digits=5}
\usepackage[version=4]{mhchem}
\usepackage{xcolor}
\usepackage{tikz}
\usepackage[most]{tcolorbox}
\usepackage{enumitem}
\usepackage{array}
\usepackage{fancyhdr}
\usepackage[hidelinks]{hyperref}
\usetikzlibrary{arrows.meta,positioning,calc}
\usepackage[a4paper,margin=2.1cm,headheight=26pt,headsep=10pt]{geometry}
\usepackage{titlesec}

\definecolor{primary}{RGB}{150,98,162}
\definecolor{primarydark}{RGB}{92,52,110}

\tcbset{boxbase/.style={enhanced,breakable,boxrule=0.9pt,arc=2.5pt,left=3mm,right=3mm,top=2.3mm,bottom=2.3mm,fonttitle=\bfseries,coltitle=white,toptitle=0.6mm,bottomtitle=0.6mm}}
\newtcolorbox[auto counter]{exercice}[1][]{boxbase,colback=primary!4,colframe=primary,title={\textsc{Exercice}~\thetcbcounter\ \ifx\relax#1\relax\relax\else\textmd{--- \itshape #1}\fi}}

\newcommand{\dS}{\Delta S}
\newcommand{\dH}{\Delta H}

\pagestyle{fancy}
\fancyhf{}
\fancyhead[L]{\small\textcolor{primarydark}{\textbf{Fiche 8} — Thème 4 : Entropie et deuxième principe}}
\fancyhead[R]{\small\textcolor{primarydark}{\textsc{Difficile}}}
\fancyfoot[C]{\small\thepage}
\renewcommand{\headrulewidth}{0.8pt}
\renewcommand{\headrule}{\hbox to\headwidth{\color{primary}\leaders\hrule height \headrulewidth\hfill}}
\setlength{\parindent}{0pt}\setlength{\parskip}{3pt}

\begin{document}

\begin{center}
\begin{tikzpicture}
\node[rectangle,rounded corners=7pt,fill=primary,text=white,minimum width=16cm,
      minimum height=1.1cm,align=center,inner sep=6pt]
  {{\large\bfseries Thème 4 — Entropie et deuxième principe}\\[1pt]{\small Niveau \textsc{Difficile} — énoncés}};
\end{tikzpicture}
\end{center}

\begin{exercice}[comparer l'ampleur du désordre]
On considère les cinq transformations :
\begin{center}\small
\begin{tabular}{@{}ll@{}}
(i) & $2\,\ce{H2O2}(l) \rightarrow 2\,\ce{H2O}(l) + \ce{O2}(g)$ \\
(ii) & $2\,\ce{NO2}(g) \rightarrow \ce{N2O4}(g)$ \\
(iii) & $\ce{I2}(s) \rightarrow \ce{I2}(g)$ \\
(iv) & $\ce{CO2}(g) \rightarrow \ce{CO2}(s)$ \\
(v) & $\ce{NaCl}(s) \rightarrow \ce{Na+}(aq) + \ce{Cl-}(aq)$ \\
\end{tabular}
\end{center}
\begin{enumerate}[label=\alph*),leftmargin=2em,topsep=2pt,itemsep=2pt]
  \item Donner le signe de $\dS$ pour chacune.
  \item Parmi celles qui \emph{augmentent} le désordre, laquelle l'augmente le plus ? Justifier.
  \item Parmi celles qui \emph{diminuent} le désordre, laquelle le diminue le plus ? Justifier.
\end{enumerate}
\end{exercice}

\begin{exercice}[la poche de froid : endothermique mais spontanée]
Une poche de froid instantané contient du nitrate d'ammonium solide qui se dissout dans l'eau :
\;$\ce{NH4NO3}(s) \rightarrow \ce{NH4+}(aq) + \ce{NO3-}(aq)$.\; Lorsqu'on l'active, la poche
\textbf{refroidit} nettement.
\begin{enumerate}[label=\alph*),leftmargin=2em,topsep=2pt,itemsep=2pt]
  \item La dissolution absorbe-t-elle ou dégage-t-elle de la chaleur ? En déduire le signe de $\dH$.
  \item Quel est le signe de l'entropie du système $\dS_{\text{système}}$ ?
  \item La dissolution se produit pourtant \emph{spontanément}. En invoquant le deuxième principe
        ($\dS_{\text{univers}}>0$), expliquer comment c'est possible malgré $\dH>0$.
  \item Si l'on ne regardait que l'\emph{énergie} ($\dH$), que prédirait-on à tort ? Que faut-il en conclure
        sur ce qui gouverne le sens d'évolution ?
\end{enumerate}
\end{exercice}

\begin{exercice}[sublimation de la neige carbonique et principes]
La neige carbonique (dioxyde de carbone solide) se sublime à l'air ambiant :
\;$\ce{CO2}(s) \rightarrow \ce{CO2}(g)$.
\begin{enumerate}[label=\alph*),leftmargin=2em,topsep=2pt,itemsep=2pt]
  \item Quel est le signe de $\dS_{\text{système}}$ ?
  \item La sublimation étant spontanée à température ambiante, que peut-on dire de $\dS_{\text{univers}}$ ?
  \item Comparer le désordre du $\ce{CO2}$ solide et du $\ce{CO2}$ gazeux.
  \item Rappeler le \textbf{troisième principe} : que vaut l'entropie du cristal parfait de $\ce{CO2}$ à
        $T=\SI{0}{\kelvin}$ ? En quoi cela distingue-t-il l'entropie de l'enthalpie (dont on ne connaît
        que les variations) ?
\end{enumerate}
\end{exercice}

\end{document}
